\documentclass[11pt]{article}

\usepackage{booktabs}
\usepackage[margin=15pt,font=small,labelfont={bf,sf},justification=justified]{caption}
\usepackage[superscript,biblabel]{cite}
\usepackage{color}
\usepackage{float}
\usepackage{graphicx}
\usepackage{lscape}
\usepackage[bbgreekl]{mathbbol}
\usepackage{mathtools}
\usepackage{multibib}
\usepackage{multirow}
\usepackage{tabularx}
\usepackage{titleref}
\usepackage{url}
\usepackage{lipsum}
\usepackage{bm}
\usepackage{tikz}
\usepackage{subfig}
\usepackage{caption}

% setup font:

\newif\ifbembo
\newif\ifcharter
\newif\iferewhon
\newif\iflibertine
\newif\iflibertinealt
\newif\ifpalantino
\newif\iftimesnewroman

% set at most one of these to be true:
\bembofalse
\chartertrue
\erewhonfalse
\libertinefalse
\libertinealtfalse
\palantinofalse
\timesnewromanfalse
% note: if all false, then this will use default latex fonts

% bembo font settings:
\ifbembo
  \usepackage[p,osf]{ETbb} % osf in text, tabular lining figures in math
  \usepackage[scaled=.95,type1]{cabin} % sans serif in style of Gill Sans
  \usepackage[varqu,varl]{zi4}% inconsolata typewriter
  \usepackage[T1]{fontenc}
  \usepackage[libertine,vvarbb]{newtxmath}
  \usepackage[cal=boondoxo,bb=boondox]{mathalfa}
\fi

% charter font settings:
\ifcharter
  \usepackage[scaled=.98]{XCharter}
  \usepackage[scaled=1.04,varqu,varl]{inconsolata}
  \usepackage[type1]{cabin}
  \usepackage[uprightscript,xcharter,vvarbb,scaled=1.05]{newtxmath}
  \usepackage[cal=euler,scr=boondoxo,bb=boondox]{mathalpha}  % prefer bb=bboldxLight when available
  \linespread{1.04}
\fi

% erewhon (utopia) font settings:
\iferewhon
  \usepackage[p,osf,scaled=.98,space]{erewhon} % scaling by .98 not really necessary
  \usepackage[varqu,varl]{inconsolata} % typewriter
  \usepackage[type1,scaled=.95]{cabin} % sans serif like Gill Sans
  \usepackage[utopia,vvarbb]{newtxmath}
\fi

% libertine font settings:
\iflibertine
  \usepackage[sb]{libertinus}
  \usepackage[T1]{fontenc}
  \usepackage{textcomp}
  \usepackage[varqu,varl]{zi4}% inconsolata for mono, not LibertineMono
  \usepackage{libertinust1math} % slanted integrals, by default
  \usepackage[scr=boondoxo,bb=boondox]{mathalpha} %Omit bb=boondox for default libertinus bb
\fi

% libertine font settings (alt):
\iflibertinealt
  \usepackage[sb]{libertinus}
  \usepackage[T1]{fontenc}
  \usepackage{textcomp}
  \usepackage{libertinust1math} % slanted integrals, by default
  \usepackage[scr=boondoxo,bb=boondox]{mathalpha} %Omit bb=boondox for default libertinus bb
\fi

% palantino font settings:
\ifpalantino
  \usepackage[largesc]{newpxtext}
  \usepackage[vvarbb]{newpxmath}
\fi

% times new roman settings:
\iftimesnewroman
  \usepackage{newtxtext} %T1 is default encoding
  \usepackage[scaled=0.95]{inconsolata} % typewriter
  \usepackage[vvarbb]{newtxmath} % vvarbb gives STIX Bbb
\fi

\usepackage[T1]{fontenc}
\usepackage[final,protrusion=true,expansion=true]{microtype}

% use sans-serif fonts for titles and all headings:
\usepackage{titling}
\usepackage{authblk} % note: titling needs to be loaded before authblk
% make titling elements \sffamily
\pretitle{\begin{center}\bfseries\sffamily\LARGE}
\posttitle{\end{center}}
\usepackage{abstract}
% make abstract title \sffamily
\renewcommand\abstractnamefont{\bfseries\sffamily\large}

%\usepackage[bf,sf]{titlesec}
\usepackage{sectsty} % note: the main alternative package, titlesec, does not work with titleref
\allsectionsfont{\sffamily}

% setup margins:
\usepackage[left=1in,right=1in,top=1in,bottom=1in,includefoot,heightrounded]{geometry}
\usepackage[hidelinks]{hyperref}

% this can be used to separate out methods or supplemental information references:
\newcites{supp}{SUPPLEMENTAL REFERENCES}

\makeatletter
\patchcmd{\LS@rot}{90}{-90}{}{}
\patchcmd{\endlandscape}{90}{-90}{}{}
\makeatother

\newcommand{\beginsupplement}{%
        \setcounter{table}{0}
        \renewcommand{\thetable}{S\arabic{table}}%
        \setcounter{figure}{0}
        \renewcommand{\thefigure}{S\arabic{figure}}%
        \setcounter{section}{0}
        \renewcommand{\thesection}{S\arabic{section}}
        \setcounter{equation}{0}
        \renewcommand{\theequation}{S\arabic{equation}}
}


\usepackage{algorithm}
\usepackage{algorithmicx}
\usepackage{algpseudocode}

\usepackage{color}
\usepackage{xspace}

%% helper macros
\newcommand{\TODO}[1]{{\color{red} \textsc{Todo} #1}}

\definecolor{vargreen}{rgb}{0.0, 0.5, 0.0}
\newcommand{\DRW}[1]{{\color{vargreen} DRW: #1}}

\newcommand{\WING}{WING\xspace}

% IBFE stuff
\newcommand{\euleriandx}{{\ensuremath{\Delta x}}}
\newcommand{\soliddom}{{\Omega^{\text{s}}_t}}
\newcommand{\soliddomO}{{\Omega^{\text{s}}_0}}

\newcommand{\Chib}{\boldsymbol{\chi}}
\newcommand{\Fb}{\boldsymbol{F}}
\newcommand{\Ub}{\boldsymbol{U}}
\newcommand{\ub}{\boldsymbol{u}}
\newcommand{\Xb}{\boldsymbol{X}}
\newcommand{\psib}{\boldsymbol{\psi}}
\newcommand{\dXb}{\, \mathrm{d}\Xb}
\newcommand{\PP}{\mathbb{P}}
\newcommand{\Sigmab}{\mathbb{S}}
\newcommand{\xb}{\boldsymbol{x}}
\newcommand{\half}{{\ensuremath{\frac{1}{2}}}}
\newcommand{\xface}{{i + \half, j, k}}
\newcommand{\yface}{{i, j + \half, k}}
\newcommand{\zface}{{i, j, k + \half}}

%% AMS band aids
\newtheorem{assumption}{Assumption}
\newtheorem{lemma}{Lemma}
\newtheorem{theorem}{Theorem}
\newtheorem{remark}{Remark}


% Set title and author information:
\title{Everything But Solvers: Optimizing Internal Finite Element Data Structures}
\author[1]{Martin Kronbichler}
\author[2,*]{David Wells}
\affil[1]{TODO}
\affil[2]{TODO}
\affil[*]{Corresponding author, \texttt{drwells@email.unc.edu}}

\begin{document}
\maketitle

\begin{abstract}
\noindent
There is much ado about fast linear solvers these days.
This is well and good.
Such solvers depend on descriptions of the geometry, parallel data layouts, and a ``long tail'' of other relevant features.
This work describes how deal.II~\cite{dealII96} modernized and adapted its core internal data structures to support high-performance solvers, mixed triangulations, and modern computer architectures.

Not a lot of people talk about efficient implementations of mixed meshes (it is unclear to me if mfem and dolfinx even support mixed meshes).
libMesh supports these things but their overreliance on inheritance makes things slow (e.g., a \texttt{Hex27} stores $27$ \texttt{Node *}s, which isn't great.
We may also want to highlight how this works and interacts with our \emph{hp}-FEM system.
\end{abstract}

%% let's have an overview of the state-of-the art, beginning with some of
%% Peter's work on Triangulation (based on the fenics paper) and other relevant
%% stuff the other packages have done since 2020 or so
\section{Introduction}
\label{sec:introduction}

\section{Literature Review}
Let's communicate what we have done in comparison to how some other relevant
packages have evolved (e.g., dolfinx).
\begin{itemize}
  \item mfem~\cite{MFEM} has published stuff on their GPU features, not so much
    on internal data structures. It does talk about some non-GPU features, but
    the focus is on solving problems with GPUs. Several of their miniapps use
    CPUs (e.g., Hooke for automatic differentiation) which is more relevant to
    this paper than their work on partial assembly etc.

  \item I also found a paper on GridapSolvers (implemented in Julia) but it
    doesn't talk about their internals.

  \item Not all apps require GPUs and large parallelism: HAZniCS
    (\url{https://arxiv.org/pdf/2210.13274}) has ``These problems do not require
    parallel computing, but they may benefit significantly from advanced
    algorithms.''

  \item The DOLFINx~\cite{DOLFINx} paper is relevant (``The design of DOLFINx
    differs fundamentally from its predecssor DOLFIN in following a
    data-oriented and functional design''). That paper says they did those
    things but doesn't provide many details (i.e., it answers `what' but not
    `how'). Relevant things:
    \begin{itemize}
      \item ``With the NBX algorithm, DOLFINx does not use any MPI all-to-all
        communication functions''

      \item Extensive use of neighborhood collectives (MPI-3). Implies that for
        very nice grids they can bound the number of neighboring processes and
        significantly simplify communication

      \item Mixed meshes were under development when this was published (2024)

      \item Not unlike our Triangulation internals (in which each object stores
        the indices of objects which bound it): they store everything as a
        graphs (adjacency lists)

      \item Focus on flat data structures makes interop between C++ and Python a
        lot easier
    \end{itemize}

  \item \emph{TODO}: find some DUNE papers. I read one on Bryne
    (\url{https://arxiv.org/pdf/2509.02378}) which talks about things they built
    on DUNE, which implements rapid prototyping on top of an FE kernel. I
    suspect a theme common to DUNE and DOLFINx is ``flat data structures enable
    interoperability with python prototypes''

    \begin{itemize}
      \item ``ALUGrid is an adaptive, loadbalancing, unstructured implementation
        of the DUNE grid interface supporting either simplices or cubes.'' (from
        \url{https://www.dune-project.org/modules/dune-alugrid/}), so no mixed
        mesh support! \texttt{SimplexGrid} and \texttt{CubeGrid} are different
        classes.

      \item Their 2015 paper has some high-level details but doesn't talk about
        low-level data structure choices (it does have some numbers on memory
        usage per vertex etc.)
    \end{itemize}
\end{itemize}

We might have to read some source code to better understand what these other
packages currently do. People don't publish these kinds of internal details
(which I suspect are generally interesting, especially if more people want to
support mixed meshes).

\section{New Optimizations Enabled by Modern C++ Standards}
\label{sec:new-optimizations}

\subsection{ReferenceCell and GeometryInfo}
deal.II 9.3~\cite{dealII93} Included preliminary support for computations on
simplex and mixed (i.e., including multiple element types) triangulations. This
was a significant milestone for deal.II, as the library historically only
supported triangulations comprised of hypercube (i.e., quadrilateral and
hexahedral) elements. While most of the foundational work was present in the 9.3
release, the developers and broader user-contributor community have continued
adding new features

\begin{enumerate}
  \item Compile-time programming is completely different now vs. 15 years ago

  \item \texttt{constexpr ReferenceCell<dim>} is about as good as
    \texttt{GeometryInfo<dim>}

  \item \texttt{constexpr} and \texttt{consteval}: do we ever need \texttt{consteval}?
\end{enumerate}

\subsection{\texttt{inplace\_vector} vs. \texttt{array} vs \texttt{vector}}
\begin{enumerate}
  \item I saw a slowdown when replacing, in the part of \texttt{tria.cc} which
    sets orientations on child faces (circa line 7910), array with inplace
    vector. Why? Is there some weird aliasing?

  \item Describe deal.II's system for generic element types (we should give
    Peter credit for this), e.g., we have \texttt{Hexahedron} and not
    \texttt{Hexahedron8} and \texttt{Hexahedron27}
\end{enumerate}

\section{Optimizing Triangulation}
\begin{enumerate}
  \item Work on flattening data structures (e.g., more inplace vector) (how much
    did it actually help?). I noticed some weirdness in Triangulation when I
    tried to replace \texttt{std::array} with
    \texttt{std\_cxx26::inplace\_vector} (it slowed down measurably in release
    mode, even though both should be contiguous stack buffers?) Maybe we could
    compare all three (i.e., include \texttt{std::vector}) possibilities and
    measure the differences (i.e., flatter data structures should help, but how
    much do they help?)

  \item Work on shrinking data structures (store one ReferenceCell per
    Triangulation instead of one per cell, variable strides) (we need to be
    careful with \texttt{std::variant}: GCC generates a vtable while clang
    creates a switch statement. Maybe this is fixed in newer versions of GCC).
    Either way we can compare a branch vs. a read

  \item For 9.8, Dominik and I rewrote a lot of the inner stride data structures
    of Triangulation (e.g., we store $N$ faces per cell, which is $4$ for
    triangulations containing only tetrahedra and $6$ for triangulations with
    hexahedra). There are a few things we can investigate there:

    \begin{itemize}
      \item Need for compile-time constants for division (PR 19944). It's an odd
        consequence of supporting mixed triangulations.

      \item Is there any difference between \emph{multiplying} by a stride known
        at compile time vs. a runtime-only constant? (see the new internal
        \texttt{TriangulationImplementation::Strides} class)
    \end{itemize}

\end{enumerate}

\section{Optimizing DoFHandler}
\begin{enumerate}
  \item Peter should also get a lot of credit here

  \item DoF lookup could be better optimized (maybe template on hp or not-hp,
    and further template on \texttt{ReferenceCell})

  \item Be careful with our use of \texttt{extern} and \texttt{inline} (maybe
    too in-the-weeds for this paper, but that was a weird effect)
\end{enumerate}

\section{Case Study in Flattening Data Structures: \texttt{KellyErrorEstimator}}
It might be worth doing some benchmarking to demonstrate the efficacy of
flattening out everything (e.g., PRs 19794, 19803).
\begin{enumerate}
  \item Describe changes (mostly switching from \texttt{std::map} to
    \texttt{Table<3, Number>})

  \item Document memory allocation changes

  \item Document callgrind or other performance benchmarks (nice learning
    experience for David: can I get \texttt{prof} to tabulate data on cache
    misses before and after? Some tooling independent of vtune would be nice)
\end{enumerate}

% main text bibliography:
\renewcommand\refname{REFERENCES}
\bibliographystyle{plain}
\bibliography{paper}

\end{document}
